\chapter{Introduzione al Data Mining e Processo KDD}
\label{ch:intro-data-mining}

\section{Definizione e Motivazioni Fondamentali}

Il costante progresso tecnologico nella generazione e memorizzazione dei dati ha portato a una crescita esponenziale delle informazioni disponibili, sia nel contesto commerciale sia in quello scientifico. Il paradigma moderno della raccolta dati si riassume spesso nel principio di raccogliere qualsiasi dato possibile, ovunque e in qualsiasi momento, nell'aspettativa che tali dati possano acquisire valore non soltanto per lo scopo originario per cui sono stati raccolti, ma anche per scopi secondari non previsti inizialmente.

\begin{definition}[Data Mining]
Il \textbf{Data Mining} è l'applicazione di tecniche ed algoritmi efficienti all'analisi di collezioni di dati di grandi dimensioni, finalizzata all'estrazione di pattern validi, inediti, potenzialmente utili e interpretabili (conoscenza nascosta o \textit{hidden knowledge}).
\end{definition}

Il Data Mining non consiste nella semplice esecuzione di query su basi di dati tradizionali (ad esempio, il recupero del saldo medio tramite SQL), bensì nella scoperta automatica o semi-automatica di strutture, dipendenze o regolarità non banali nei dati.

\subsection{Prospettive Applicative}
Le motivazioni alla base dell'adozione delle tecniche di Data Mining possono essere distinte principalmente secondo due punti di vista:

\begin{itemize}
    \item \textbf{Punto di vista commerciale:} Le aziende gestiscono volumi enormi di dati (\textit{data warehousing}): dati web (Yahoo, Facebook), transazioni di e-commerce e vendite al dettaglio (Amazon), registrazioni di carte di credito e conti correnti bancari. A fronte di una potenza di calcolo sempre più economica e accessibile, la pressione competitiva richiede di estrarre valore da tali moli di dati per migliorare e personalizzare i servizi offerti (ad esempio nel \textit{Customer Relationship Management} - CRM).
    \item \textbf{Punto di vista scientifico e sociale:} In discipline quali l'astronomia (\textit{Sky Survey Data}), l'osservazione della Terra (i satelliti del sistema NASA EOSDIS raccolgono petabyte di dati ambientali all'anno), la genomica (espressione genica, proteomica), la neuroimaging (fMRI) e le simulazioni numeriche complesse, i dati vengono generati a velocità tali da rendere impossibile l'analisi manuale. Il Data Mining supporta gli scienziati nell'analisi automatica e nella formulazione guidata di nuove ipotesi.
\end{itemize}

I Big Data rappresentano a tutti gli effetti un proxy della vita sociale (comportamenti d'acquisto, stili di vita, opinioni, reti di relazioni, mobilità sul territorio), offrendo grandi opportunità per affrontare problematiche globali: dall'efficientamento della sanità alla ricerca di fonti di energia pulita, dalla previsione degli impatti climatici all'ottimizzazione della produzione agricola.

\subsection{Una Confluenza di Discipline}
Il Data Mining non è una disciplina isolata, bensì una confluenza sinergica di diverse aree dell'informatica e della matematica applicata:

\begin{center}
\begin{tikzpicture}[
    node distance=2.2cm,
    every node/.style={font=\small},
    center_node/.style={circle, draw=blue!70, fill=blue!10, thick, inner sep=6pt, align=center, font=\bfseries\small},
    satellite/.style={rectangle, rounded corners, draw=gray!70, fill=gray!5, thick, inner sep=5pt, align=center}
]
    \node[center_node] (dm) {Data\\Mining};
    \node[satellite, above left=1.2cm and 1.8cm of dm] (stat) {Statistica};
    \node[satellite, above right=1.2cm and 1.8cm of dm] (db) {Sistemi di\\Basi di Dati};
    \node[satellite, below left=1.2cm and 1.8cm of dm] (ml) {Machine\\Learning};
    \node[satellite, below right=1.2cm and 1.8cm of dm] (alg) {Teoria degli\\Algoritmi};
    \node[satellite, above=1.5cm of dm] (vis) {Visualizzazione};
    \node[satellite, below=1.5cm of dm] (other) {Altre Discipline\\(HPC, Ottimizzazione)};

    \draw[thick, -Stealth, color=blue!60] (stat) -- (dm);
    \draw[thick, -Stealth, color=blue!60] (db) -- (dm);
    \draw[thick, -Stealth, color=blue!60] (ml) -- (dm);
    \draw[thick, -Stealth, color=blue!60] (alg) -- (dm);
    \draw[thick, -Stealth, color=blue!60] (vis) -- (dm);
    \draw[thick, -Stealth, color=blue!60] (other) -- (dm);
\end{tikzpicture}
\end{center}

\section{Tipologie di Origine del Dato: Primari e Secondari}
Nell'impostazione di un progetto analitico, è fondamentale distinguere l'origine dei dati analizzati:
\begin{itemize}
    \item \textbf{Dati Primari (\textit{Primary Data}):} Dati originali rilevati specificamente per un preciso obiettivo di indagine. Non hanno subito alterazioni o manipolazioni preventive da parte di terzi.
    \item \textbf{Dati Secondari (\textit{Secondary Data}):} Dati che erano già stati raccolti in precedenza per scopi operativi o differenti da quello corrente e resi disponibili all'analista. Spesso derivano dalla fusione di sorgenti disomogenee e necessitano di profonde operazioni di normalizzazione e verifica della qualità.
\end{itemize}

\section{Il Processo KDD (Knowledge Discovery in Databases)}
Il Data Mining costituisce una fase specifica, sebbene centrale, all'interno del più ampio processo di scoperta della conoscenza, noto come \textbf{KDD} (\textit{Knowledge Discovery in Databases}). Il KDD è un processo iterativo e interattivo composto da diverse fasi sequenziali con cicli di retroazione:

\begin{center}
\begin{tikzpicture}[
    node distance=1.1cm,
    font=\small,
    box/.style={rectangle, rounded corners, draw=blue!80, fill=blue!5, thick, minimum width=4.5cm, minimum height=0.8cm, align=center},
    data/.style={cylinder, shape border rotate=90, draw=gray!80, fill=gray!10, aspect=0.25, minimum height=0.9cm, minimum width=3.2cm, align=center}
]
    \node[data] (db) {Sorgenti Eterogenee\\(Basi di dati, Log, Web)};
    \node[box, below=0.8cm of db] (clean) {1. Data Cleaning \& Integration};
    \node[data, below=0.8cm of clean] (dw) {Data Warehouse};
    \node[box, below=0.8cm of dw] (select) {2. Data Selection \& Transformation};
    \node[data, below=0.8cm of select] (taskdata) {Dati Rilevanti (Task-relevant)};
    \node[box, below=0.8cm of taskdata, draw=red!80, fill=red!5] (dm) {3. Data Mining (Algoritmi)};
    \node[box, below=0.8cm of dm] (eval) {4. Pattern Evaluation \&\\Knowledge Presentation};

    \draw[-Stealth, thick] (db) -- (clean);
    \draw[-Stealth, thick] (clean) -- (dw);
    \draw[-Stealth, thick] (dw) -- (select);
    \draw[-Stealth, thick] (select) -- (taskdata);
    \draw[-Stealth, thick] (taskdata) -- (dm);
    \draw[-Stealth, thick] (dm) -- (eval);

    \draw[-Stealth, dashed, thick, color=gray!80] (eval.east) to [bend right=45] node[right, font=\scriptsize] {Raffinamento iterativo} (select.east);
\end{tikzpicture}
\end{center}

Ciascuna fase assolve compiti specifici:
\begin{enumerate}
    \item \textbf{Business Understanding:} Comprensione degli obiettivi di dominio e traduzione del problema di business in un problema formale di analisi dati.
    \item \textbf{Data Cleaning e Data Integration:} 
    \begin{itemize}
        \item \textit{Data Cleaning}: rimozione di rumore, correzione di errori sintattici e gestione delle incoerenze.
        \item \textit{Data Integration}: aggregazione e riconciliazione di dati provenienti da fonti multiple ed eterogenee, confluendo tipicamente in un \textbf{Data Warehouse} (un repository strutturato e ottimizzato per l'analisi e il supporto alle decisioni).
    \end{itemize}
    \item \textbf{Data Selection e Data Transformation:}
    \begin{itemize}
        \item \textit{Data Selection}: estrazione del sottoinsieme di attributi e record strettamente pertinenti al task investigato (\textit{task-relevant data}).
        \item \textit{Data Transformation}: trasformazione delle feature mediante aggregazioni, riscalamenti, normalizzazioni o discretizzazioni affinché siano adatte all'algoritmo.
    \end{itemize}
    \item \textbf{Data Mining:} Esecuzione vera e propria di algoritmi intelligenti per l'estrazione di pattern e modelli matematici.
    \item \textbf{Pattern Evaluation e Knowledge Presentation:} Identificazione dei pattern realmente significativi tramite metriche oggettive e soggettive, eliminazione della ridondanza e visualizzazione grafica dei risultati mediante strumenti interpretabili dall'utente finale.
\end{enumerate}

\section{Tassonomia dei Compiti di Data Mining}

Seguendo la classificazione di Fayyad et al. \cite{fayyad1996}, i compiti di Data Mining si dividono in due macro-categorie:
\begin{itemize}
    \item \textbf{Metodi Predittivi (\textit{Prediction Methods}):} Mirano a stimare il valore futuro o sconosciuto di una variabile target specifica a partire dai valori di altre variabili osservate (attributi predittori).
    \item \textbf{Metodi Descrittivi (\textit{Description Methods}):} Mirano a identificare pattern leggibili ed interpretabili dall'essere umano che descrivano la struttura intrinseca e le relazioni nascoste nei dati.
\end{itemize}

\subsection{Compiti Predittivi: La Classificazione}
La \textbf{classificazione} ha come obiettivo l'apprendimento di una funzione obiettivo $f$ che mappa ogni record di attributi $x$ in una specifica etichetta di classe discreta $y \in \mathcal{Y}$.

\begin{figure}[htbp]
\centering
\begin{tikzpicture}[font=\small, node distance=1.5cm]
    \node[rectangle, draw=blue!70, fill=blue!5, thick, text width=3.2cm, align=center] (train) {\textbf{Training Set}\\Record storici con classe nota $y$};
    \node[rectangle, draw=red!70, fill=red!5, thick, text width=2.8cm, align=center, right=1.2cm of train] (learn) {\textbf{Algoritmo di Apprendimento}\\(es. Albero Decisionale)};
    \node[rectangle, draw=green!70, fill=green!5, thick, text width=2.6cm, align=center, right=1.2cm of learn] (model) {\textbf{Modello Classificatore}\\$y = f(x)$};
    \node[rectangle, draw=gray!70, fill=gray!5, thick, text width=3.2cm, align=center, below=1.2cm of model] (test) {\textbf{Test Set / Dati Nuovi}\\Record con classe da stimare};
    \node[rectangle, draw=purple!70, fill=purple!5, thick, text width=2.6cm, align=center, left=1.2cm of test] (pred) {\textbf{Previsione}\\Etichetta di classe};

    \draw[-Stealth, thick] (train) -- (learn);
    \draw[-Stealth, thick] (learn) -- (model);
    \draw[-Stealth, thick] (model) -- (test);
    \draw[-Stealth, thick] (test) -- (pred);
\end{tikzpicture}
\caption{Workflow concettuale del processo di classificazione.}
\label{fig:classification-workflow}
\end{figure}

Il dataset viene suddiviso in due insiemi disgiunti:
\begin{itemize}
    \item \textbf{Training Set:} Utilizzato per addestrare il classificatore.
    \item \textbf{Test Set:} Utilizzato per valutare la capacità di generalizzazione del modello su istanze non osservate durante la fase di apprendimento.
\end{itemize}

\begin{example}[Applicazioni di Classificazione]
Tra le principali e pi\`u diffuse aree applicative dei modelli di classificazione troviamo:
\begin{itemize}
    \item \textbf{Valutazione del Merito Creditizio (\textit{Credit Worthiness}):} Predire se un richiedente prestito rimborser\`a il debito (\textit{Yes/No}) sulla base di occupazione, livello di istruzione e anni di residenza.
    \item \textbf{Rilevamento Frodi Bancarie:} Classificare le transazioni con carta di credito come lecite o fraudolente. Le feature includono l'importo, l'orario, la posizione geografica e il profilo storico delle spese.
    \item \textbf{Churn Prediction:} Stimare la probabilit\`a che un cliente di telecomunicazioni abbandoni il servizio per passare alla concorrenza, analizzando la durata delle chiamate, il volume del traffico dati e i reclami registrati.
    \item \textbf{Altre applicazioni:} Categorizzazione automatica di notizie (sport, finanza, politica), diagnosi clinica (tumore benigno o maligno), classificazione di coperture del suolo da immagini satellitari.
\end{itemize}
\end{example}

\subsection{Compiti Descrittivi}

\subsubsection{1. Clustering (Raggruppamento)}
Il \textbf{clustering} consiste nell'organizzare un insieme di oggetti in gruppi distinti (detti \textit{cluster}), in modo tale che:
\begin{enumerate}
    \item La distanza tra oggetti appartenenti allo stesso cluster (\textbf{distanza intra-cluster}) sia \textbf{minimizzata} (massima omogeneità interna).
    \item La distanza tra oggetti appartenenti a cluster differenti (\textbf{distanza inter-cluster}) sia \textbf{massimizzata} (massima separazione esterna).
\end{enumerate}

\begin{figure}[htbp]
\centering
\begin{tikzpicture}[scale=0.85]
    % Cluster 1
    \draw[thick, fill=blue!10, draw=blue!60, rounded corners=15pt] (-3.5,-1) rectangle (-1.2,1.5);
    \fill[blue!70] (-3,0.8) circle (3pt);
    \fill[blue!70] (-2.2,1.1) circle (3pt);
    \fill[blue!70] (-2.7,0.1) circle (3pt);
    \fill[blue!70] (-1.8,-0.4) circle (3pt);
    \node[blue!80!black, font=\bfseries\scriptsize] at (-2.35,-0.7) {Cluster 1};

    % Cluster 2
    \draw[thick, fill=green!10, draw=green!60, rounded corners=15pt] (1.2,-1) rectangle (3.5,1.5);
    \fill[green!60!black] (1.7,0.9) circle (3pt);
    \fill[green!60!black] (2.4,1.1) circle (3pt);
    \fill[green!60!black] (3.0,0.2) circle (3pt);
    \fill[green!60!black] (2.1,-0.5) circle (3pt);
    \node[green!60!black, font=\bfseries\scriptsize] at (2.35,-0.7) {Cluster 2};

    % Intra-cluster distance arrow
    \draw[Stealth-Stealth, thick, blue!80] (-2.7,0.1) -- node[above, font=\tiny] {intra-distanza min} (-1.8,-0.4);

    % Inter-cluster distance arrow
    \draw[Stealth-Stealth, very thick, red!80] (-1.2,0.5) -- node[above, font=\scriptsize\bfseries] {inter-distanza max} (1.2,0.5);
\end{tikzpicture}
\caption{Principio geometrico fondamentale del clustering.}
\label{fig:clustering-principle}
\end{figure}

Il clustering è una tecnica non supervisionata: non richiede classi a priori.
\begin{itemize}
    \item \textbf{Segmentazione di Mercato:} Raggruppare i clienti in base a dati socio-demografici e abitudini di acquisto per indirizzare campagne di marketing differenziate.
    \item \textbf{Clustering Documentale:} Raggruppare documenti testuali simili in base alle frequenze congiunte delle parole chiave (ad esempio, l'analisi dell'archivio email aziendale Enron).
    \item \textbf{Summarization e Scienze della Terra:} Ridurre la complessità di enormi serie spaziali partizionando le misurazioni di temperatura marina superficiale (\textit{Sea Surface Temperature} - SST) e produzione primaria (\textit{Net Primary Production} - NPP) negli emisferi settentrionale e meridionale.
\end{itemize}

\subsubsection{2. Scoperta di Regole di Associazione (Association Analysis)}
Dato un insieme di transazioni, ciascuna costituita da un sottoinsieme di elementi presi da un insieme globale di articoli (\textit{items}), l'analisi delle associazioni si prefigge di scoprire dipendenze probabilistiche nella forma di regole:
\begin{equation}
X \implies Y
\end{equation}
dove $X$ e $Y$ sono itemset disgiunti ($X \cap Y = \emptyset$). La regola indica che la presenza degli articoli nell'antecedente $X$ predice sistematicamente la contestuale presenza dell'articolo o degli articoli nel conseguente $Y$.

\begin{example}[Market Basket Analysis]
Considerando uno storico di scontrini di cassa:
\begin{table}[htbp]
\centering
\begin{tabular}{cl}
\toprule
\textbf{TID} & \textbf{Articoli Acquistati} \\
\midrule
1 & Pane, Coca-Cola, Latte \\
2 & Birra, Pane \\
3 & Birra, Coca-Cola, Pannolini, Latte \\
4 & Birra, Pane, Pannolini, Latte \\
5 & Coca-Cola, Pannolini, Latte \\
\bottomrule
\end{tabular}
\end{table}
L'algoritmo estrae regole ad elevata evidenza statistica, quali:
\[
\{\text{Latte}\} \implies \{\text{Coca-Cola}\}, \quad \{\text{Pannolini}, \text{Latte}\} \implies \{\text{Birra}\}
\]
Tali informazioni guidano la disposizione strategica dei beni sugli scaffali e la progettazione di offerte cross-selling.
\end{example}

\textbf{Altre applicazioni:} Diagnosi di allarmi nelle reti di telecomunicazione (individuazione di pattern temporali di allarmi di rete che scattano congiuntamente per identificare guasti sistemici alla radice); informatica medica (correlazione tra combinazioni di sintomi clinici, esami ematochimici e patologie emergenti).

\subsubsection{3. Rilevamento di Anomalie (Anomaly / Deviation / Change Detection)}
Consiste nell'identificare oggetti che deviano in modo significativo dal comportamento medio o atteso della popolazione. 
Le anomalie possono essere generate da guasti strumentali o errori di inserimento (e in tal caso rappresentano \textit{rumore} da epurare), oppure riflettere eventi eccezionali di estremo interesse applicativo (frodi creditizie, intrusioni informatiche contro le reti aziendali, monitoraggio di repentini disboscamenti della copertura forestale tramite sensori remoti).

\section{Sfide Metodologiche e Computazionali nel Data Mining}
L'applicazione di metodi tradizionali di analisi statistica e database è spesso inadeguata a causa di cinque sfide primarie:
\begin{enumerate}
    \item \textbf{Scalabilità (\textit{Scalability}):} Gli algoritmi devono gestire volumi di dati di ordine di terabyte o petabyte senza che la complessità computazionale e l'occupazione di memoria esplodano.
    \item \textbf{Alta Dimensionalità (\textit{High Dimensionality}):} La presenza di centinaia o migliaia di attributi espone alla cosiddetta \textit{maledizione della dimensionalità} (\textit{curse of dimensionality}), in cui lo spazio geometrico diventa estremamente vuoto e le distanze tendono a perdere discriminatività.
    \item \textbf{Dati Complessi ed Eterogenei:} Necessità di elaborare strutture semi-strutturate o non strutturate (grafi, testi, sequenze genomiche, serie spaziotemporali).
    \item \textbf{Proprietà e Distribuzione dei Dati:} I dati sono frequentemente distribuiti geograficamente tra molteplici organizzazioni e soggetti a stringenti vincoli legali di privacy e sicurezza.
    \item \textbf{Analisi Non Tradizionali:} Esigenza di estrarre conoscenza da flussi continui in tempo reale (\textit{data stream mining}) senza possibilità di memorizzare l'intero dataset.
\end{enumerate}
